% =====================================================================
\section{Introduction}
% =====================================================================
Large language models are now routinely deployed as \emph{judges}: instead of recruiting human annotators to label
data, rate outputs, or score quality, researchers prompt an LLM to do it for a fraction of the cost
\citep{zheng2023mtbench,liu2023geval,gilardi2023chatgpt}. The practice is spreading faster than the evidence that
licenses it. Two questions remain under-answered. \textbf{When} does an LLM judge agree with humans closely enough to
substitute for them, and \textbf{what} governs that agreement, model scale, task, or something about the human
labels themselves?

We answer these questions for \emph{open, locally-servable} models. This is a deliberate scope, not a limitation:
proprietary judges are expensive, non-reproducible, and change without notice, so for most researchers the realistic
question is whether an open model they can run themselves is a trustworthy judge. Our entire empirical program uses
open models served through a standard vLLM endpoint at temperature~0 with released harnesses.

One methodological commitment runs through the paper: an LLM judge must be measured against the \emph{human-human
agreement ceiling} for the task, not against a single ``gold'' label. For subjective constructs there is no single
correct answer \citep{aroyo2015truth,plank2022humanlabel}; the relevant question is whether the judge agrees with
humans as well as humans agree with each other. We compute that ceiling directly from the per-item annotator
distributions that the datasets already provide, so our human reference is real, independent, and inherited from the
data, not added by us.

\paragraph{Contributions.} (1) \textbf{Study 1:} a PRISMA systematic review of \Nstudies{} studies across
\Ndomains{} domains (every reference independently web-verified) that maps, against the human ceiling, where
LLM-as-a-judge is reliable and where it is not. (2) \textbf{Study 2:} a controlled replication showing that the most
powerful open models now meet or exceed the established proprietary-judge correlations on the verifiable aspects of a
flagship benchmark, with bootstrap confidence intervals. (3) \textbf{Study 3:} a new meta-evaluation of a six-model
open panel (9B-754B) on six datasets with human label distributions (\STHRn{} judgments), establishing that
\emph{item-level human disagreement}, not model scale or data contamination, is the binding constraint on judge
reliability. (4) A decision framework and an eight-point protocol for designing LLM-as-a-judge experiments. All
code, corpora, search logs, and harnesses are released.
